\documentclass[12pt, a4paper]{article}
\usepackage[UTF8]{ctex}
\usepackage{amsmath}
\usepackage{amssymb}
\usepackage{geometry}
\usepackage{enumitem}
\usepackage{fancyhdr}
\usepackage{titlesec}

% 页面设置
\geometry{left=2.5cm, right=2.5cm, top=2.5cm, bottom=2.5cm}

% 页眉页脚
\pagestyle{fancy}
\fancyhf{}
\rhead{专升本数据结构习题集}
\lhead{分类整理}
\cfoot{\thepage}

% 标题格式
\titleformat{\section}{\large\bfseries}{\thesection}{1em}{}
\titleformat{\subsection}{\normalsize\bfseries}{\thesubsection}{1em}{}

\title{\textbf{第五章 数组和广义表}}
\author{}
\date{}

\begin{document}

\section*{选择题}

\begin{enumerate}[label=\arabic*.]
	\item 一维数组和线性表的区别是（ ）
	      \begin{enumerate}
		      \item[A.] 前者长度固定，后者长度可变
		      \item[B.] 后者长度固定，前者长度可变
		      \item[C.] 两者长度均固定
		      \item[D.] 两者长度均可变
	      \end{enumerate}

	\item 二维数组 M 的成员是 6 个字符组成的串，行下标 i 的范围从 0 到 8，列下标 j 的范围从 1 到 10，则存放 M 至少需要（①）个字节，M 的第 8 列和第 5 行共占（②）个字节，若 M 按行优先方式存储，元素 M[8][5] 的起始地址与当 M 按列优先方式存储时元素（③）的起始地址一致。
	      \begin{enumerate}
		      \item[①] A. 90 \quad B. 180 \quad C. 240 \quad D. 540
		      \item[②] A. 108 \quad B. 114 \quad C. 54 \quad D. 60
		      \item[③] A. M[8][5] \quad B. M[3][10] \quad C. M[5][8] \quad D. M[0][9]
	      \end{enumerate}

	\item 二维数组 M 的元素是 4 个字符（每个字符占一个存储单元）组成的串，行下标 i 范围从 0 到 4，列下标 j 的范围从 0 到 5，M 按行优先存储时元素 M[3][5] 的起始地址与 M 按列优先存储时元素（ ）的起始地址相同。
	      \begin{enumerate}
		      \item[A.] M[2][4]
		      \item[B.] M[3][4]
		      \item[C.] M[4][3]
		      \item[D.] M[5][2]
	      \end{enumerate}

	\item 数组 A 中，每个元素 A 的长度为 3 个字节，行下标 i 从 1 到 8，列下标 j 从 1 到 10，从首地址 SA 开始连续存放在存储器内，存放该数组至少需要的单元数是（ ）
	      \begin{enumerate}
		      \item[A.] 80
		      \item[B.] 100
		      \item[C.] 240
		      \item[D.] 270
	      \end{enumerate}

	\item 数组 A 中每个元素的长度为 3 个字节，行下标 i 从 1 到 8，列下标 j 从 1 到 10，从首地址 SA 开始连续存放在存储器内，该数组按行存放时，元素 A[8][5] 的起始地址为（ ）
	      \begin{enumerate}
		      \item[A.] SA+141
		      \item[B.] SA+144
		      \item[C.] SA+222
		      \item[D.] SA+225
	      \end{enumerate}

	\item 数组 A 中每个元素的长度为 3 个字节，行下标 i 从 1 到 8，列下标 j 从 1 到 10，从首地址 SA 开始连续存放在存储器内，该数组按列存放时，A[5][8] 的起始地址为（ ）
	      \begin{enumerate}
		      \item[A.] SA+141
		      \item[B.] SA+180
		      \item[C.] SA+222
		      \item[D.] SA+225
	      \end{enumerate}

	\item 设有有一个 10 阶的对称矩阵 A，采用压缩存储方式，以行序为主存储，$a_{1,1}$ 为第一个元素，其存储地址为 1，每个元素占 1 个地址空间，则 $a_{8,5}$ 的地址为（ ）
	      \begin{enumerate}
		      \item[A.] 13
		      \item[B.] 33
		      \item[C.] 18
		      \item[D.] 40
	      \end{enumerate}

	\item 一个 $n \times n$ 的对称矩阵，如果以行或列为主序放入内存，则容量为（ ）
	      \begin{enumerate}
		      \item[A.] $n \times n$
		      \item[B.] $n \times n / 2$
		      \item[C.] $n \times (n+1) / 2$
		      \item[D.] $(n+1) \times (n+1) / 2$
	      \end{enumerate}

	\item 稀疏矩阵一般的压缩存储方法有两种，即（ ）
	      \begin{enumerate}
		      \item[A.] 二维数组和三维数组
		      \item[B.] 三元组和散列
		      \item[C.] 三元组和十字链表
		      \item[D.] 散列和十字链表
	      \end{enumerate}

	\item 若采用三元组压缩技术存储稀疏矩阵，只要把每个元素的行下标和列下标互换，就完成了对该矩阵的转置运算，这种观点（ ）
	      \begin{enumerate}
		      \item[A.] 正确
		      \item[B.] 错误
	      \end{enumerate}

	\item 数组 $A[0..5, 0..6]$ 的每个元素占五个字节，将其按列优先次序存储在起始地址为 1000 的内存单元中，则元素 A[5, 5] 的地址是（ ）
	      \begin{enumerate}
		      \item[A.] 1175
		      \item[B.] 1180
		      \item[C.] 1205
		      \item[D.] 1210
	      \end{enumerate}

	\item 将一个 $A[1..100, 1..100]$ 的三对角矩阵，按行优先存入一维数组 $B[1..298]$ 中，A 中元素 $A_{66,65}$ (即该元素下标 $i=66, j=65$)，在 B 数组中的位置 K 为（ ）
	      \begin{enumerate}
		      \item[A.] 198
		      \item[B.] 195
		      \item[C.] 197
		      \item[D.] 201
	      \end{enumerate}

	\item 有一个 $100 \times 90$ 的稀疏矩阵，非 0 元素有 10 个，设每个整型数据占 2 个字节，则用三元组表示该矩阵时，所需的字节数是（ ）
	      \begin{enumerate}
		      \item[A.] 60
		      \item[B.] 66
		      \item[C.] 18000
		      \item[D.] 33
	      \end{enumerate}

	\item 数组 $A[0..4, -1..-3, 5..7]$ 中含有元素的个数（ ）
	      \begin{enumerate}
		      \item[A.] 55
		      \item[B.] 45
		      \item[C.] 36
		      \item[D.] 16
	      \end{enumerate}

	\item 对数组经常进行的操作有（ ）
	      \begin{enumerate}
		      \item[A.] 建立与删除
		      \item[B.] 索引与修改
		      \item[C.] 查找与修改
		      \item[D.] 查找与索引
	      \end{enumerate}

	\item 下面说法不正确的是（ ）
	      \begin{enumerate}
		      \item[A.] 广义表的表头总是一个广义表
		      \item[B.] 广义表的表尾总是一个广义表
		      \item[C.] 广义表难以用顺序存储结构
		      \item[D.] 广义表可以是一个多层次的结构
	      \end{enumerate}

	\item 广义表 $((a), a)$ 的表头是（ ），表尾是（ ）
	      \begin{enumerate}
		      \item[A.] a
		      \item[B.] b
		      \item[C.] (a)
		      \item[D.] ((a))
	      \end{enumerate}

	\item 广义表 $((a))$ 的表头是（ ），表尾是（ ）
	      \begin{enumerate}
		      \item[A.] a
		      \item[B.] (a)
		      \item[C.] ()
		      \item[D.] ((a))
	      \end{enumerate}

	\item 广义表 $((a,b),c,d)$ 的表头是（ ），表尾是（ ）
	      \begin{enumerate}
		      \item[A.] a
		      \item[B.] b
		      \item[C.] (a, b)
		      \item[D.] (c, d)
	      \end{enumerate}

	\item 广义表 $(a,b,c,d)$ 的表头是（ ），表尾是（ ）
	      \begin{enumerate}
		      \item[A.] a
		      \item[B.] b
		      \item[C.] (a, b)
		      \item[D.] (b, c, d)
	      \end{enumerate}

	\item 一个广义表的表头总是一个广义表，这个断言是（ ）
	      \begin{enumerate}
		      \item[A.] 正确的
		      \item[B.] 错误的
	      \end{enumerate}

	\item 一个广义表的表尾总是一个广义表，这个断言是（ ）
	      \begin{enumerate}
		      \item[A.] 正确的
		      \item[B.] 错误的
	      \end{enumerate}

	\item 已知广义表 $L=((x, y, z), (u, t, w))$，从 L 表中取出原子 t 的运算是（ ）
	      \begin{enumerate}
		      \item[A.] head(tail(tail(L)))
		      \item[B.] tail(head(head(tail(L))))
		      \item[C.] head(tail(head(tail(L))))
		      \item[D.] head(head(tail(tail(L))))
	      \end{enumerate}

	\item 若广义表 $A=(a,b,(c,d),(e,(f,g)))$，则下面式子的值为（ ）：Head(tail(head(tail(tail(A)))))
	      \begin{enumerate}
		      \item[A.] (g)
		      \item[B.] (d)
		      \item[C.] (c)
		      \item[D.] d
	      \end{enumerate}

	\item 若广义表 A 满足 Head(A)=Tail(A)，则 A 为（ ）
	      \begin{enumerate}
		      \item[A.] ()
		      \item[B.] (())
		      \item[C.] ((),())
		      \item[D.] ((),(),())
	      \end{enumerate}

	\item 一个向量第一个元素的存储地址是 100，每个元素的长度为 2，则第 5 个元素的地址是（ ）
	      \begin{enumerate}
		      \item[A.] 110
		      \item[B.] 108
		      \item[C.] 100
		      \item[D.] 120
	      \end{enumerate}

	\item 数据在计算机存储器内表示时，物理地址与逻辑地址相同并且是连续的，称之为（ ）
	      \begin{enumerate}
		      \item[A.] 存储结构
		      \item[B.] 逻辑结构
		      \item[C.] 顺序存储结构
		      \item[D.] 链式存储结构
	      \end{enumerate}

	\item 对矩阵进行压缩存储是为了（ ）
	      \begin{enumerate}
		      \item[A.] 方便运算
		      \item[B.] 方便存储
		      \item[C.] 提高运算速度
		      \item[D.] 减少存储空间
	      \end{enumerate}

	\item 数组 A 中，每个元素的长度为 3 个字节，行下标 i 从 1 到 8，列下标 j 从 1 到 10，从首地址 SA 开始连续存放的存储器内，该数组按行存放，元素 A[8][5] 的起始地址为（ ）
	      \begin{enumerate}
		      \item[A.] SA+141
		      \item[B.] SA+144
		      \item[C.] SA+222
		      \item[D.] SA+225
	      \end{enumerate}

	\item 数组 A 中，每个元素的长度为 3 个字节，行下标 i 从 1 到 8，列下标 j 从 1 到 10，从首地址 SA 开始连续存放的存储器内，该数组按行存放，元素 A[8][5] 的起始地址为（ ）
	      \begin{enumerate}
		      \item[A.] SA+141
		      \item[B.] SA+144
		      \item[C.] SA+222
		      \item[D.] SA+225
	      \end{enumerate}

	\item 设二维数组 $A[1..m, 1..n]$ 按行存储在数组 B 中，则二维数组元素 $A[i, j]$ 在一维数组 B 中的下标为（ ）
	      \begin{enumerate}
		      \item[A.] $n \times (i-1) + j$
		      \item[B.] $n \times (i-1) + j - 1$
		      \item[C.] $i \times (j-1)$
		      \item[D.] $j \times m + i - 1$
	      \end{enumerate}

	\item 若声明一个浮点数数组如下：float average[]=new float[30]；假设该数组的内存起始位置为 200，average[15] 的内存地址是（ ）
	      \begin{enumerate}
		      \item[A.] 214
		      \item[B.] 215
		      \item[C.] 260
		      \item[D.] 256
	      \end{enumerate}

	\item 字符数组 $a[1..100]$ 采用顺序存储，$a[6]$ 地址是 517，那么 a 的首地址为（ ）
	      \begin{enumerate}
		      \item[A.] 510
		      \item[B.] 512
		      \item[C.] 514
		      \item[D.] 516
	      \end{enumerate}

	\item 数组 a[1..256] 采用顺序存储，a 的首地址为 10，每个元素占 2 字节，则 a[21] 的地址是（ ）
	      \begin{enumerate}
		      \item[A.] 10
		      \item[B.] 70
		      \item[C.] 50
		      \item[D.] 30
	      \end{enumerate}

	\item 数组 a[1..m] 采用顺序存储，a[1] 和 a[m] 地址分别为 1024 和 1150，每个元素占 2 字节，则 m 是（ ）
	      \begin{enumerate}
		      \item[A.] 63
		      \item[B.] 64
		      \item[C.] 65
		      \item[D.] 66
	      \end{enumerate}

	\item 设有一个 10 阶的对称矩阵 A，采用压缩存储方式，以行序为主存储，$a_{1,1}$ 为第一个元素，其存储地址为 1，每个元素占 1 个地址空间，则 $a_{8,5}$ 的地址为（ ）
	      \begin{enumerate}
		      \item[A.] 13
		      \item[B.] 33
		      \item[C.] 18
		      \item[D.] 40
	      \end{enumerate}

	\item 下面关于 B 树和 B+ 树的叙述中，不正确的结论是（ ）
	      \begin{enumerate}
		      \item[A.] B 树和 B+ 树都能有效的支持顺序查找
		      \item[B.] B 树和 B+ 树都能有效的支持随机查找
		      \item[C.] B 树和 B+ 树都是平衡的多叉树
		      \item[D.] B 树和 B+ 树都可用于文件索引结构
	      \end{enumerate}

	\item 下面叙述错误的是（ ）
	      \begin{enumerate}
		      \item[A.] 对矩阵进行压缩存储后无法实现对其元素进行随机访问
		      \item[B.] 空串的长度为零
		      \item[C.] 借助于栈可以实现对图的深度优先遍历
	      \end{enumerate}

	\item 对称矩阵压缩存储时，需要存储的元素个数为（ ）
	      \begin{enumerate}
		      \item[A.] $n \times n$
		      \item[B.] $n \times (n+1) / 2$
		      \item[C.] $n \times (n-1) / 2$
		      \item[D.] $n \times n / 2$
	      \end{enumerate}

	\item 下面关于压缩存储的叙述中，正确的是（ ）
	      \begin{enumerate}
		      \item[A.] 稀疏矩阵必须采用三元组存储
		      \item[B.] 对称矩阵可以采用压缩存储
		      \item[C.] 三角矩阵不能采用压缩存储
		      \item[D.] 稀疏矩阵不能用数组存储
	      \end{enumerate}
\end{enumerate}

\section*{填空题}

\begin{enumerate}
	\item 当广义表中的每个元素都是原子（ ）时，广义表便成了（ ）。
	\item 广义表 $(a, (a,b), d, e, ((i, j), k))$ 的长度是（ ），深度是（ ）。
	\item 一维数组的存储结构是（ ），逻辑结构是（ ），对二维数组来说分（ ）和（ ）两种不同的存储方式。
	\item 已知二维数组采用行为主方式存储，每个元素占 k 个存储单元，并且第一个元素的存储地址是 LOC(A[0][0])，则 A[i][j] 的地址是（ ）。
	\item 二维数组 A[10][20] 采用以列为主方式存储，每个元素占一个存储单元，并且 A[0][0] 的存储地址是 200，则 A[6][12] 的地址是（ ）。
	\item 二维数组 A[10..20][5..10] 采用行为主方式存储，每个元素占四个存储单元，并且 A[10][5] 的存储地址是 1000，则 A[18][9] 的地址是（ ）。
	\item 有一个 10 阶对称矩阵 A，采用压缩存储方式 (以行序为主存储且 A[0][0]=1)，则 A[8][5] 的地址是（ ）。
	\item 设 n 行 n 列的下三角矩阵 A 已压缩到一维数组 $S[1..n*(n+1)/2+1]$ 中，若按行序为主存储，当 $i \ge j$ 时，A[i][j] 对应的 S 中的存储位置是（ ）。
	\item 一个对称矩阵，采用压缩存储方式存储其容量为（ ），若不采用压缩存储则容量为（ ）；三角矩阵压缩存储时其容量为（ ）。
	\item 二维数组 M 的每个元素是 6 个字符组成的串，行下标 i 的范围从 0 到 8，列下标 j 的范围从 1 到 10，则存放 M 至少需要 \underline{540} 个字节；M 的第 8 列和第 5 行共占 \underline{108} 个字节。
	\item 稀疏矩阵一般的压缩存储方法有两种，即 \underline{三元组表} 和 \underline{十字链表}。
	\item 广义表 $((a), ((b), c), (((d))))$ 的长度是 \underline{3}，深度是 \underline{4}。
	\item 一维数组的存储结构是顺序存储，逻辑结构是线性结构，对二维数组来说分行优先和列优先两种不同的存储方式。
	\item 已知二维数组采用行为主方式存储，每个元素占 $k$ 个存储单元，并且第一个元素的存储地址是 $\text{LOC}(A[0][0])$，则 $A[i][j]$ 的地址是 $\text{LOC}(A[0][0]) + (i \times n + j) \times k$。

	\item 二维数组 $A[10][20]$ 采用以列为主方式存储，每个元素占一个存储单元，并且 $A[0][0]$ 的存储地址是 $200$，则 $A[6][12]$ 的地址是 $200 + 12 \times 10 + 6 = 326$。

	\item 二维数组 $A[10..20][5..10]$ 采用行为主方式存储，每个元素占四个存储单元，并且 $A[10][5]$ 的存储地址是 $1000$，则 $A[18][9]$ 的地址是 $1000 + ((18-10) \times 6 + (9-5)) \times 4 = 1000 + (8 \times 6 + 4) \times 4 = 1000 + 52 \times 4 = 1208$。

	\item 有一个 10 阶对称矩阵 $A$，采用压缩存储方式（以行序为主存储且 $A[0][0]=1$），则 $A[8][5]$ 的地址是 $33$。
	\item 设 n 行 n 列的下三角矩阵 A 已压缩到一维数组 $S[1..n*(n+1)/2+1]$ 中，若按行序为主存储，当 $i \ge j$ 时，A[i][j] 对应的 S 中的存储位置是 i(i-1)/2 + j。
	\item 一个对称矩阵，采用压缩存储方式存储其容量为 $n(n+1)/2$，若不采用压缩存储则容量为 $n^2$；三角矩阵压缩存储时其容量为 $n(n+1)/2$。
\end{enumerate}

\section*{判断题}

\begin{enumerate}
	\item 若一个广义表的表头为空表，则此广义表也为空表。（ ）
	\item 广义表是线性表的推广，是一类线性结构。（ ）
	\item 一个稀疏矩阵采用三元组形式表示，若把三元组中行下标和列下标的值互换，则完成了矩阵的转置运算。（ ）
	\item 数组是一种复杂的数据结构，数组元素之间的关系既不是线性的，也不是树型的。（ ）
	\item 线性表可以看成是广义表的特例，如果广义表中的每个元素都是原子，则广义表便成为线性表。（ ）
	\item 广义表中原子的个数即为广义表的长度。（ ）
	\item 稀疏矩阵压缩存储后，必会失去随机存取功能。（ ）
	\item 数组元素的下标值越大，存取时间越长。（ ）
	\item 用邻接矩阵法存储一个图时，在不考虑压缩存储的情况下，所占用的存储空间大小只与图中结点个数有关，而与图的边数无关。（ ）
	\item 一个广义表的表头总是一个广义表。（ ）
	\item 一个广义表的表尾总是一个广义表。（ ）
	\item 广义表 $(((a), b), c)$ 的表头是 $((a), b)$，表尾是 $(c)$。（ ）
	\item 数组是一种复杂的数据结构。（ ）
	\item 稀疏矩阵一般采用压缩存储。（ ）
	\item 广义表可以递归定义。（ ）
	\item 二维数组是其数据元素为线性表的线性表。（ ）
\end{enumerate}

\section*{应用题}

\begin{enumerate}
	\item 已知广义表 $L=((((a))), ((b)), (c), d)$，试利用 head 和 tail 运算把原子项 c 从 L 中分离出来。
	\item 画出下列广义表的两种存储结构图：$(O, A, (B, (C, D)), (E, F))$。
	\item 已知广义表 $A=(((b,c),d),(a),((a),((b,c),d)),e,())$
	      \begin{enumerate}
		      \item 画出其一种存储结构图；
		      \item 写出表的长度与深度，表头与表尾。
	      \end{enumerate}
	\item 稀疏矩阵 A 如图所示，画出以下各种表示法：（1）行优先三元组表示法；（2）列优先三元组表示法；（3）伪地址表示法；（4）十字链表示法。
	      \[
		      \begin{pmatrix}
			      15 & 0  & 0 & 22 & 0  & -15 \\
			      0  & 13 & 3 & 0  & 0  & 0   \\
			      0  & 0  & 0 & -6 & 0  & 0   \\
			      0  & 0  & 0 & 0  & 91 & 0   \\
			      0  & 0  & 0 & 0  & 0  & 0   \\
			      0  & 0  & 2 & 8  & 0  & 0
		      \end{pmatrix}
	      \]
\end{enumerate}

\end{document}
